\begin{claim}{C.61}{Proved}{Compact prefix families prevent fiber intervals from returning to an open orbit}
Retain the original based plane loop $f_t$, its
generated group $G$, and the closed set
\[
 Y=\R^2\setminus\bigcup\{O:O\text{ is an open }G\text{-orbit}\}.
\]
For each open orbit under consideration, choose one
actual global quotient, its original universal
source cover $p_O:X_O\to O$, and one real target
lift $F_O:X_O\to\R$ as in C.43--C.44.
Fix these choices on each orbit.
Write $\widetilde f_t^O$ for the based source lift
of $f_t|_O$.

Suppose $x_n\in O_n$ converges to $x\in Y$, and
$s_n,t_n\in[0,1]$. Choose any source lift
$\widetilde x_n$ of each $x_n$ and require the
exact real equality
\[
 F_{O_n}(\widetilde f_{s_n}^{O_n}\widetilde x_n)
   =F_{O_n}(\widetilde f_{t_n}^{O_n}\widetilde x_n).
                                                               \tag{C.61.1}
\]
Let $I_n$ be the unique compact interval in that
proper real fiber joining those two lifted
endpoints, and put $J_n=p_{O_n}(I_n)$.
Then every compact set $K\subset\R^2\setminus Y$
is eventually disjoint from $J_n$:
\[
                  J_n\cap K=\varnothing
                      \quad\hbox{for all sufficiently large }n.
                                                               \tag{C.61.2}
\]
In particular every finite ambient accumulation
point of these interval images lies in $Y$.
The statement allows arbitrary variation of the
open orbit, either source covering type, and the
original target deck labels. It requires equality
in (C.61.1), rather than equality modulo a circle
target's deck group.
\end{claim}

\begin{proof}
C.43 gives connected proper real-line fibers and
an actual product chart
$X_O\cong\R^2$ in which $F_O(u,v)=v$.
Every based source lift preserves this partition:
its real quotient equation has the form
$F_O\widetilde f_t^O=b_tF_O$ with $b_t$ increasing.
The lift commutes with each original source deck
map $\delta$. Indeed,
$\widetilde f_t^O\delta(\widetilde f_t^O)^{-1}\delta^{-1}$
is a continuous deck-valued path starting at the
identity, and hence is the identity.
Also C.44 retains the exact equation
\[
                 F_O\delta=F_O+\eta_O(\delta),           \tag{C.61.3}
\]
where the label is zero for a line target and is
an integer for the specified circle cover.
Thus a deck map carries the interval in (C.61.1)
to the corresponding interval in its translated
real fiber, with no change of endpoint ambiguity.

For any one fixed $O$, the actual family
\[
 \mathcal P_O=
 \{\widetilde f_t^O(\widetilde f_s^O)^{-1}:
                                      s,t\in[0,1]\}
                                                               \tag{C.61.4}
\]
is compact in $\operatorname{Homeo}(X_O)$, with
inverse convergence. This follows from continuity
of the original lifted path and multiplication,
on the compact parameter square.
In the product chart its elements have the form
\[
            g(u,v)=(h_g(u,v),b_g(v)),                    \tag{C.61.5}
\]
where $h_g(\,\cdot\,,v)$ is increasing and onto.
For a convergent sequence $g_n\to g$ in
$\mathcal P_O$ and $v_n\to v$, these fiber maps
are uniformly proper along $v_n$. Explicitly,
given $M$, choose $A$ with $h_g(A,v)>M+1$.
Then $h_{g_n}(A,v_n)>M$ for all sufficiently
large $n$, and monotonicity gives
$h_{g_n}(u_n,v_n)>M$ whenever $u_n>A$.
The lower bound follows by choosing $A$ with
$h_g(A,v)<-M-1$.
Consequently
\[
 \begin{split}
 u_n\longrightarrow+\infty
   &\ \Longrightarrow\ h_{g_n}(u_n,v_n)\longrightarrow+\infty,\\
 u_n\longrightarrow-\infty
   &\ \Longrightarrow\ h_{g_n}(u_n,v_n)\longrightarrow-\infty.
 \end{split}                                             \tag{C.61.6}
\]
These estimates use the whole compact family
of lifted homeomorphisms and their actual
increasing fiber maps.

Suppose (C.61.2) fails. After a subsequence there
are $y_n\in J_n\cap K$ with $y_n\to y\in K$.
The point $y$ lies in some open orbit $O$.
Eventually $y_n\in O$, while $y_n\in O_n$.
Distinct $G$-orbits are disjoint, so $O_n=O$
on this tail. All source covers and the chosen
real quotient are now the same fixed maps.
Pass to $s_n\to s$ and $t_n\to t$.

Choose $\widetilde y_n\in I_n$ over $y_n$.
A covering neighborhood at $y$ supplies lifts
$\widehat y_n\to\widehat y$ in one compact source
chart. Choose the actual deck map $\delta_n$
with $\delta_n\widetilde y_n=\widehat y_n$.
Put
\[
 \begin{split}
 a_n&=\delta_n\widetilde f_{s_n}^O\widetilde x_n,\\
 c_n&=\delta_n\widetilde f_{t_n}^O\widetilde x_n,\qquad
 g_n=\widetilde f_{t_n}^O(\widetilde f_{s_n}^O)^{-1}.
 \end{split}                                             \tag{C.61.7}
\]
Commutation with the very same deck map
$\delta_n$ gives the exact equality $c_n=g_na_n$.
By (C.61.3), the interval $\delta_n I_n$ joins
$a_n$ to $c_n$ in one real fiber and contains
$\widehat y_n$.
Write
$a_n=(u_n,v_n)$ and
$\widehat y_n=(w_n,v_n)$.
Then $v_n\to v$ and $w_n\to w$ are finite.
Equation (C.61.1) also gives $b_{g_n}(v_n)=v_n$.

The first endpoints escape every compact subset
of $X_O$. Indeed,
\[
                 p_O(a_n)=f_{s_n}(x_n)\longrightarrow f_s(x).
 \]
The set $Y$ is $G$-invariant, so $f_s(x)\in Y$.
A convergent source subsequence would instead
project to a point of $O$.
Since $v_n$ is bounded, this proves
$|u_n|\to\infty$. Select either the positive
or the negative infinite subsequence.
Compactness of (C.61.4) gives $g_n\to g$,
and (C.61.6) forces the first coordinate of
$c_n=g_na_n$ to escape to the same end.
The real interval between these two coordinates
cannot contain the bounded $w_n$.
This contradicts
$\widehat y_n\in\delta_n I_n$ and proves
(C.61.2).

No uniform choice of charts on different
orbits was made: an alleged open accumulation
point forces a single actual orbit before
any chart is used.
The result does not bound accumulation along
$Y$, nor exclude escape to ambient infinity.
In particular it does not yet give the required
control near one specified nonopen orbit or
a lift of the nontrivial quotient contraction
left open after C.60.
\end{proof}
